\documentclass[11pt]{article}
\usepackage[margin=1in]{geometry}
\usepackage{amsmath,amssymb,amsthm,enumitem}
\usepackage{hyperref}

\newtheorem{theorem}{Theorem}
\newtheorem{lemma}{Lemma}
\newtheorem{corollary}{Corollary}
\newtheorem{definition}{Definition}
\newtheorem{remark}{Remark}
\newtheorem{example}{Example}

\newcommand{\adj}{\operatorname{adj}}
\newcommand{\spann}{\operatorname{span}}
\newcommand{\Ran}{\operatorname{Ran}}
\newcommand{\kry}{\mathcal K}
\newcommand{\nobs}{\mathcal N_{\mathrm{obs}}}
\newcommand{\Vresp}{\mathcal V_{\mathrm{resp}}}
\newcommand{\nudiss}{\nu_{\mathrm{diss}}}

\title{Three Central Theorems of the Sector-Resolved Response Classification:\\
Exact Structural No-Go, Dissipative Asymptotic Hierarchy,\\
and Protected Channels under Singular Dissipation}
\author{}
\date{}

\begin{document}
\maketitle

\begin{abstract}
We prove the three central theorems that upgrade the binary EIT no-go criterion
to a three-tier classification of sector-mediated coherent response in
finite-dimensional Markovian weak-probe systems: (I) an exact structural no-go
theorem in Kalman controllability/observability form; (II) a dissipative
asymptotic hierarchy with a uniform-in-frequency expansion and a well-defined
suppression index $\nudiss$; (III) a protected-channel theorem for singular,
generally non-Hermitian, dissipation via Riesz projections, together with its
EIT sector-cut corollary. We then prove that the induced classification
$\nu\in\{\infty\}\cup(0,\infty)\cup\{0\}$ is exhaustive and mutually exclusive,
that all judged quantities are basis independent, and that the associated
decision algorithm terminates in finitely many steps. Building blocks proved in
the companion document \texttt{eit\_nogo\_proofs.tex} (Theorems 2A/2B, 3, 4, 5,
6 there) are cited where reused.
\end{abstract}

% ==================================================================
\section{Framework, standing assumptions, and basis independence}
% ==================================================================

\subsection{Arena}

Throughout, $\mathcal V\cong\mathbb C^{N}$ is a finite-dimensional (Liouville)
space. The weak-probe linear response of the open system is encoded in a
generator family
\begin{equation}
A_\Gamma(z)=\Gamma D+B(z),
\label{eq:AGamma}
\end{equation}
where $\Gamma\ge0$ is an overall dissipation scale, $D$ is the ($z$-independent)
dissipation matrix, and $B(z)$ collects coherent evolution, detunings, and
fixed-rate dissipators; $B$ is a matrix-valued function, analytic on an open
frequency domain $\Omega\subset\mathbb C$ and \emph{bounded on compact subsets}
of $\Omega$. The scalar transfer function of interest is
\begin{equation}
F_\Gamma(z)=p^\dagger A_\Gamma(z)^{-1}c ,
\label{eq:transfer}
\end{equation}
with fixed source vector $c$ (probe injection) and readout vector $p$
(detected coherence).

For the \emph{sector-cut} response we use the realization established in the
companion document (Theorem 2B there): for a designated sector $S$,
\begin{equation}
\delta\chi_S(z)
 =\chi_{\mathrm{full}}(z)-\chi_{\mathrm{cut}}^{(S)}(z)
 =\ell^\dagger (zI-M)^{-1} r ,
\label{eq:realization}
\end{equation}
a finite-dimensional state-space realization with system matrix $M$ acting on
a realization space of dimension $n\le N^2$, input vector $r$, and output
vector $\ell$. Unless stated otherwise, $M$, $r$, $\ell$ are $z$-independent
(the standard weak-probe realization used throughout the repository); Remark
\ref{rem:zdep} discusses the $z$-dependent case.

\subsection{Standing assumptions}

\begin{enumerate}[label=(A\arabic*),nosep]
\item \emph{Finite dimension}: all spaces are finite dimensional.
\item \emph{Markovianity}: the dynamics is generated by a (possibly
      probe-dressed) Lindbladian; $D$ and $B(z)$ arise from its
      block decomposition.
\item \emph{Weak probe}: response is linear in the probe; all statements
      concern first order in the probe field.
\item \emph{Passivity}: no gain; $\Re\operatorname{spec}(-A_\Gamma)\le0$ on the
      physical domain, so resolvents exist off the spectrum.
\end{enumerate}
No claim is made outside this scope (strong-probe nonlinearity, non-Markovian
baths, continua, ultrastrong coupling, time-dependent control, and gain media
are excluded).

\subsection{Basis independence of the judged quantities}

\begin{lemma}[Similarity invariance]
\label{lem:basis}
Let $T$ be invertible and transform the realization \eqref{eq:realization} by
$(M,r,\ell)\mapsto(TMT^{-1},\,Tr,\,T^{-\dagger}\ell)$. Then
\begin{enumerate}[label=(\roman*),nosep]
\item the transfer function $\delta\chi_S(z)$ is unchanged for every $z$;
\item every moment $\mu_k=\ell^\dagger M^k r$ is unchanged;
\item for the family \eqref{eq:AGamma} transformed by
      $(D,B,c,p)\mapsto(TDT^{-1},TBT^{-1},Tc,T^{-\dagger}p)$, the function
      $F_\Gamma(z)$, hence every coefficient of any asymptotic expansion of
      $F_\Gamma$ in $\Gamma^{-1}$ (in particular $\nudiss$, $F_0$, $F_1$, and
      $\delta\chi_{S,0}$ defined below), is unchanged.
\end{enumerate}
\end{lemma}

\begin{proof}
(i) $(T^{-\dagger}\ell)^\dagger(zI-TMT^{-1})^{-1}Tr
=\ell^\dagger T^{-1}\,T(zI-M)^{-1}T^{-1}\,Tr=\ell^\dagger(zI-M)^{-1}r$.
(ii) Identically, $\ell^\dagger T^{-1}(TMT^{-1})^kTr=\ell^\dagger M^kr$.
(iii) The same conjugation collapses inside
$p^\dagger A_\Gamma^{-1}c$, so $F_\Gamma(z)$ is invariant as a function of
$(\Gamma,z)$. Asymptotic coefficients of an invariant function are invariant:
if $F_\Gamma=\sum_{k\le m}a_k\Gamma^{-k}+O(\Gamma^{-m-1})$ in two bases, the
difference of the two polynomial parts is $O(\Gamma^{-m-1})$ and hence zero,
coefficient by coefficient.
\end{proof}

\begin{lemma}[Realization independence]
\label{lem:kalman}
$\delta\chi_S(z)$ depends only on the restriction of $(M,r,\ell)$ to the
\emph{minimal} (Kalman) subsystem: the quotient of the reachable subspace
$\kry(M,r)=\spann\{r,Mr,\dots,M^{n-1}r\}$ by its intersection with the
unobservable subspace $\nobs(M,\ell)=\bigcap_{k\ge0}\ker(\ell^\dagger M^k)$.
Consequently all quantities of Lemma \ref{lem:basis} may be computed in any
realization of the same transfer function.
\end{lemma}

\begin{proof}
$\kry(M,r)$ is $M$-invariant and contains $r$; $\nobs(M,\ell)$ is
$M$-invariant and is annihilated by $\ell^\dagger M^k$ for all $k$. Expanding
$(zI-M)^{-1}=\sum_{k\ge0}z^{-k-1}M^k$ for $|z|>\|M\|$, every Laurent
coefficient $\ell^\dagger M^kr$ is computed from the trajectory
$\{M^kr\}\subset\kry(M,r)$, and any component of that trajectory inside
$\nobs(M,\ell)$ contributes zero. Two realizations with equal transfer
functions have equal Laurent coefficients (uniqueness of Laurent expansions),
i.e.\ equal moment sequences, and every quantity above is a function of the
moment sequence and of the transfer function itself.
\end{proof}

% ==================================================================
\section{Theorem I: sector-resolved transfer-function no-go}
% ==================================================================

\begin{theorem}[I: Exact Structural No-Go]
\label{thm:I}
Let $\delta\chi_S(z)=\ell^\dagger(zI-M)^{-1}r$ with $\dim=n$ and
$z$-independent $(M,r,\ell)$. The following are equivalent:
\begin{enumerate}[label=(\roman*),nosep]
\item $\mu_k:=\ell^\dagger M^k r=0$ for $k=0,1,\dots,n-1$;
\item $\kry(M,r)\subseteq\nobs(M,\ell)$;
\item $\ell^\dagger q(M)\,r=0$ for every polynomial $q$;
\item $\delta\chi_S(z)=0$ for every $z\notin\operatorname{spec}M$
      (exact zero at all frequencies).
\end{enumerate}
A system satisfying these conditions is in \textbf{Class I}, and we set
$\nu=\infty$.
\end{theorem}

\begin{proof}
\emph{(i)$\Rightarrow$(iii).} By Cayley--Hamilton, $M^n=\sum_{k<n}c_kM^k$ for
scalars $c_k$, so $\mu_n=\sum_{k<n}c_k\mu_k=0$; multiplying the
Cayley--Hamilton relation repeatedly by $M$ and inducting gives $\mu_k=0$ for
all $k\ge0$, hence $\ell^\dagger q(M)r=0$ for every polynomial $q$.

\emph{(iii)$\Rightarrow$(ii).} $\kry(M,r)$ is spanned by
$\{M^kr\}_{k<n}$. For any $v=M^kr$ and any $j\ge0$,
$\ell^\dagger M^jv=\ell^\dagger M^{j+k}r=0$ by (iii); hence
$v\in\nobs(M,\ell)$, and $\nobs$ contains the whole span.

\emph{(ii)$\Rightarrow$(i).} $r\in\kry(M,r)\subseteq\nobs$ and $\nobs$ is
$M$-invariant, so $M^kr\in\nobs$ and in particular
$\mu_k=\ell^\dagger M^kr=0$ for all $k$.

\emph{(iii)$\Rightarrow$(iv).} For $z\notin\operatorname{spec}M$,
$(zI-M)^{-1}=\adj(zI-M)/\det(zI-M)$, and $\adj(zI-M)$ is a polynomial in $M$
with $z$-dependent scalar coefficients: with characteristic polynomial
$\chi(t)=t^n+a_{n-1}t^{n-1}+\dots+a_0$ ($a_n:=1$), the matrix
\[
B(z):=\sum_{k=0}^{n-1}\Bigl(\sum_{j=k+1}^{n}a_j\,z^{\,j-k-1}\Bigr)M^{k}
\]
satisfies $(zI-M)B(z)=\chi(z)I$ (telescoping plus $\chi(M)=0$), hence
$B(z)=\adj(zI-M)$ by uniqueness of the inverse. Then
$\delta\chi_S(z)=\ell^\dagger B(z)r/\chi(z)=0$ by (iii), for \emph{every} $z$
off the spectrum --- no large-$|z|$ restriction.

\emph{(iv)$\Rightarrow$(i).} For $|z|>\|M\|$ expand
$\delta\chi_S(z)=\sum_{k\ge0}\mu_k z^{-k-1}$ (Neumann series). A convergent
Laurent series vanishing on an open annulus has all coefficients zero, so
$\mu_k=0$ for all $k$, in particular $k<n$.
\end{proof}

\begin{remark}[Finite certificate]
Condition (i) is a \emph{finite} test: $n$ scalar evaluations certify an exact
all-frequency zero. This is Step 1 of the decision algorithm
(Section~\ref{sec:algorithm}).
\end{remark}

\begin{corollary}[Mechanism checklist]
\label{cor:mechanisms}
Each of the following implies the conditions of Theorem \ref{thm:I}:
\begin{enumerate}[label=(\alph*),nosep]
\item \emph{Symmetry separation}: orthogonal projectors $\{P_\alpha\}$ with
      $\sum_\alpha P_\alpha=I$, $[M,P_\alpha]=0$, $P_\beta r=r$,
      $\ell^\dagger P_\alpha=\ell^\dagger$, $\alpha\ne\beta$
      (companion Theorem 5);
\item \emph{Unreachability}: $r$ lies in an $M$-invariant subspace $\mathcal W$
      with $\ell\perp\mathcal W$;
\item \emph{Unobservability}: $\ell^\dagger$ annihilates an $M$-invariant
      subspace containing $r$ (this is (ii) verbatim);
\item \emph{All-order cancellation}: $\mu_k=0$, $k<n$, without any invariant
      subspace visible in the physical basis (``accidental'' zeros).
\end{enumerate}
\end{corollary}

\begin{proof}
(a) With $\alpha\neq\beta$: $M^kr=M^kP_\beta r=P_\beta M^kr$, so
$\ell^\dagger M^kr=\ell^\dagger P_\alpha P_\beta M^kr=0$.
(b) $M^kr\in\mathcal W$ for all $k$, so $\mu_k=\ell^\dagger M^kr=0$.
(c) Is (ii). (d) Is (i). \end{proof}

\begin{remark}[$z$-dependent realizations]
\label{rem:zdep}
If $M=M(z)$ through residual $z$-dependence of $B$, Theorem \ref{thm:I}
holds at each fixed $z$ (equivalence of (i)--(iii) with $M=M(z)$, and
$\delta\chi_S(z)=0$ at that $z$ by the adjugate identity, whose polynomial
degree bound is uniform). The all-frequency statement (iv) then requires the
fixed-$z$ conditions to hold on the whole domain; when the conditions come
from a $z$-independent mechanism of Corollary \ref{cor:mechanisms}(a)--(c),
this is automatic.
\end{remark}

% ==================================================================
\section{Theorem II: dissipative asymptotic hierarchy}
% ==================================================================

We now return to the family \eqref{eq:AGamma}. Define the \emph{response
subspace} $\Vresp\subseteq\mathcal V$ as the smallest subspace containing $c$
that is invariant under $D$ and under $B(z)$ for all $z\in\Omega$, enlarged by
the corresponding adjoint-invariant subspace generated by $p$; by Lemmas
\ref{lem:basis}--\ref{lem:kalman} we may and do restrict all operators to
$\Vresp$ without changing $F_\Gamma$, and we write $N$ for its dimension.

\begin{theorem}[II: Dissipative Asymptotic Hierarchy]
\label{thm:II}
Assume $D$ is invertible on $\Vresp$. Define
\[
\mu_k(z)=p^\dagger\bigl[D^{-1}B(z)\bigr]^kD^{-1}c ,\qquad k\ge0 .
\]
Let $K\subset\Omega$ be compact and $X_K:=\sup_{z\in K}\|D^{-1}B(z)\|<\infty$.
Then:
\begin{enumerate}[label=(\roman*),nosep]
\item for every $\Gamma>X_K$ and $z\in K$,
\[
F_\Gamma(z)=\sum_{k=0}^{\infty}(-1)^k\,\Gamma^{-(k+1)}\mu_k(z),
\]
absolutely convergent, with the uniform tail bound
\[
\Bigl|F_\Gamma(z)-\sum_{k=0}^{m}(-1)^k\Gamma^{-(k+1)}\mu_k(z)\Bigr|
\le\frac{\|p\|\,\|D^{-1}\|\,\|c\|}{\Gamma}\cdot
\frac{(X_K/\Gamma)^{m+1}}{1-X_K/\Gamma}
\qquad\text{for all }z\in K;
\]
\item if $\mu_0\equiv\dots\equiv\mu_{m-1}\equiv0$ on $K$ and
      $\mu_m\not\equiv0$, then
\[
F_\Gamma(z)=(-1)^m\mu_m(z)\,\Gamma^{-(m+1)}+O(\Gamma^{-(m+2)})
\qquad\text{uniformly for }z\in K;
\]
\item the first index $m$ with $\mu_m\not\equiv0$, if it exists, satisfies
      $m\le N-1$; otherwise all $\mu_k\equiv0$ and $F_\Gamma\equiv0$
      (Class I). The \emph{dissipative suppression index}
\[
\boxed{\;\nudiss:=m+1\;}
\]
is well defined, basis independent, and takes values in
$\{1,\dots,N\}\cup\{\infty\}$.
\end{enumerate}
Systems with $0<\nudiss<\infty$ form \textbf{Class II}, with
$\nu=\nudiss$.
\end{theorem}

\begin{proof}
\emph{(i).} $X_K<\infty$ because $z\mapsto\|D^{-1}B(z)\|$ is continuous on the
compact $K$ ($B$ analytic, hence continuous). Factor
$A_\Gamma=\Gamma D+B=D\bigl(\Gamma I+D^{-1}B\bigr)$, so
\[
F_\Gamma(z)=p^\dagger\bigl(\Gamma I+D^{-1}B(z)\bigr)^{-1}D^{-1}c
=\frac1\Gamma\,p^\dagger\Bigl(I+\tfrac{D^{-1}B(z)}{\Gamma}\Bigr)^{-1}D^{-1}c .
\]
For $\Gamma>X_K$ the Neumann series
$(I+Y)^{-1}=\sum_{k\ge0}(-Y)^k$ with $Y=D^{-1}B(z)/\Gamma$,
$\|Y\|\le X_K/\Gamma<1$, converges absolutely and uniformly on $K$;
term-by-term contraction with $p^\dagger$ and $D^{-1}c$ gives the series and
the geometric tail bound, whose constants do not depend on $z\in K$.

\emph{(ii).} Truncate (i) at order $m$: the surviving leading term is
$(-1)^m\mu_m(z)\Gamma^{-(m+1)}$ and the tail is bounded by
$\mathrm{const}\cdot\Gamma^{-1}(X_K/\Gamma)^{m+1}/(1-X_K/\Gamma)
=O(\Gamma^{-(m+2)})$ with a $z$-uniform constant on $K$.

\emph{(iii).} Fix $z_0\in K$ with (eventually) $\mu_m(z_0)\ne0$. The scalars
$\mu_k(z_0)$ are moments of the pair $\bigl(X(z_0),\nu\bigr)$ with
$X=-D^{-1}B(z_0)$, $\nu=D^{-1}c$ against $p$; by Theorem \ref{thm:I}
(equivalence (i)$\Leftrightarrow$(iii) applied on $\Vresp$, $\dim=N$), if
$\mu_k(z_0)=0$ for all $k\le N-1$ then $\mu_k(z_0)=0$ for \emph{all} $k$.
Hence either the first nonzero index obeys $m\le N-1$ at some $z_0\in K$, or
all moments vanish identically on $K$, in which case (i) forces
$F_\Gamma\equiv0$ on $K$ for all large $\Gamma$, and, $F_\Gamma$ being
rational in $\Gamma$, for all $\Gamma$ off a finite set: Class I. Basis
independence of $m$ (hence $\nudiss$) is Lemma \ref{lem:basis}(iii).
\end{proof}

\begin{corollary}[Sector-graph selection rule]
\label{cor:graph}
Suppose $\Vresp$ splits into sectors preserved by $D$ and by the diagonal part
of $B=B_{\mathrm{diag}}+V$, and let $d$ be the minimal number of
$V$-insertions connecting the sector of $c$ to the sector of $p$. Then
$\mu_k\equiv0$ for $k<d$, hence $\nudiss\ge d+1$: the suppression index counts
the effective path length through the sector graph.
\end{corollary}

\begin{proof}
Each $\mu_k$ expands into words in $D^{-1}B_{\mathrm{diag}}$ and $D^{-1}V$;
a word with fewer than $d$ factors of $V$ cannot map the sector of $c$ to that
of $p$, so its overlap with $p$ vanishes (companion document, remark after
Theorem 3, reproved here in the $B$-decomposition).
\end{proof}

% ==================================================================
\section{Theorem III: protected channels under singular dissipation}
% ==================================================================

When $D$ is singular the hierarchy of Theorem \ref{thm:II} does not apply on
$\ker D$; instead an $O(1)$ channel may survive. In Liouville space $D$ is in
general \emph{not} Hermitian, so we work with Riesz (spectral) projections
rather than orthogonal ones.

\begin{lemma}[Riesz splitting of a semisimple kernel]
\label{lem:riesz}
Let $D$ have $0$ as a semisimple eigenvalue (no Jordan block), or be
invertible ($P=0$ below). Let
\[
P=\frac1{2\pi i}\oint_{\gamma}(\zeta I-D)^{-1}\,d\zeta ,
\]
with $\gamma$ a small circle around $0$ enclosing no other eigenvalue, and
$Q=I-P$. Then $P^2=P$, $PD=DP=0$, $\mathcal V=P\mathcal V\oplus Q\mathcal V$
(direct, generally non-orthogonal sum), $\Ran P=\ker D$, and
$D_Q:=D|_{Q\mathcal V}$ is invertible on $Q\mathcal V$.
\end{lemma}

\begin{proof}
$P$ is the standard Riesz projection: idempotency and the direct-sum
decomposition into the spectral subspaces inside/outside $\gamma$ are the
holomorphic functional calculus facts for matrices, and $[P,D]=0$ because $P$
is a contour integral of resolvents of $D$. The spectral subspace of the
eigenvalue $0$ is the generalized eigenspace $\ker D^{N}$; semisimplicity
means $\ker D^{N}=\ker D$, hence $\Ran P=\ker D$ and $DP=PD=0$. On
$Q\mathcal V$ the spectrum of $D$ is $\operatorname{spec}D\setminus\{0\}$,
bounded away from $0$, so $D_Q$ is invertible.
\end{proof}

\begin{theorem}[III: Protected Channel under Singular Dissipation]
\label{thm:III}
Assume, on a compact frequency window $K\subset\Omega$:
\begin{enumerate}[label=(H\arabic*),nosep]
\item $0$ is a semisimple eigenvalue of $D$; $P,Q$ as in Lemma
      \ref{lem:riesz};
\item $D_Q=QDQ|_{Q\mathcal V}$ is invertible (automatic from Lemma
      \ref{lem:riesz});
\item the \emph{protected block} $B_P(z)=PB(z)P|_{P\mathcal V}$ is invertible
      for every $z\in K$;
\item fixed-coupling regime: $\sup_{z\in K}\|B(z)\|<\infty$
      (no $\Gamma$-dependence of $B$).
\end{enumerate}
Then there is $\Gamma_*<\infty$ such that for $\Gamma>\Gamma_*$ and $z\in K$,
\[
F_\Gamma(z)=F_0(z)+\Gamma^{-1}F_1(z)+O(\Gamma^{-2})
\qquad\text{uniformly on }K,
\]
with
\begin{align}
F_0(z)&=p^\dagger P\,B_P(z)^{-1}\,Pc ,
\label{eq:F0}\\
F_1(z)&=p^\dagger\Bigl[
 P B_P^{-1}\,PBQ\,D_Q^{-1}\,QBP\,B_P^{-1}P
 \;-\;PB_P^{-1}\,PBQ\,D_Q^{-1}Q
 \;-\;QD_Q^{-1}\,QBP\,B_P^{-1}P
 \;+\;QD_Q^{-1}Q
\Bigr]c .
\label{eq:F1}
\end{align}
If $F_0\not\equiv0$ on $K$, then $F_\Gamma=O(1)$: a finite response survives
arbitrarily strong dissipation. Such systems form \textbf{Class III}
($\nu=0$).
\end{theorem}

\begin{proof}
Work in the (oblique) decomposition $\mathcal V=P\mathcal V\oplus Q\mathcal V$
of Lemma \ref{lem:riesz}; since $PD=DP=0$,
\[
A_\Gamma=\Gamma D+B=
\begin{pmatrix}B_{PP}&B_{PQ}\\ B_{QP}&\Gamma D_Q+B_{QQ}\end{pmatrix},
\qquad
B_{XY}:=XB(z)Y \ \ (X,Y\in\{P,Q\}),
\]
where $B_{PP}=B_P$. The obliquity of the splitting only enters through the
finite constants $\|P\|,\|Q\|$, which are $\Gamma$- and $z$-independent.

\emph{Lower-right block.} For
$\Gamma>\Gamma_1:=\|D_Q^{-1}\|\sup_K\|B\|\,\|Q\|^2$,
\[
(\Gamma D_Q+B_{QQ})^{-1}
=\frac1\Gamma D_Q^{-1}\Bigl(I+\tfrac{D_Q^{-1}B_{QQ}}{\Gamma}\Bigr)^{-1}
=\frac1\Gamma D_Q^{-1}-\frac1{\Gamma^2}D_Q^{-1}B_{QQ}D_Q^{-1}
 +O(\Gamma^{-3}),
\]
by the Neumann series, uniformly on $K$ by (H4).

\emph{Schur complement on the protected block.}
\[
S_P(\Gamma,z)=B_P-B_{PQ}\,(\Gamma D_Q+B_{QQ})^{-1}B_{QP}
=B_P-\frac1\Gamma B_{PQ}D_Q^{-1}B_{QP}+O(\Gamma^{-2}).
\]
By (H3), $B_P^{-1}$ exists on $K$ and, $K$ compact and $B_P$ continuous,
$\sup_K\|B_P^{-1}\|<\infty$; hence for $\Gamma$ beyond some $\Gamma_*\ge\Gamma_1$
the correction is a small perturbation and
\[
S_P^{-1}=B_P^{-1}
 +\frac1\Gamma B_P^{-1}B_{PQ}D_Q^{-1}B_{QP}B_P^{-1}
 +O(\Gamma^{-2})
\qquad\text{uniformly on }K .
\]

\emph{Block solve.} Writing $c=Pc+Qc$, $p^\dagger=p^\dagger P+p^\dagger Q$ and
solving $A_\Gamma x=c$ by eliminating the $Q$ block:
\begin{align*}
x_P&=S_P^{-1}\bigl(Pc-B_{PQ}(\Gamma D_Q+B_{QQ})^{-1}Qc\bigr)
   =B_P^{-1}Pc
   +\tfrac1\Gamma\bigl[B_P^{-1}B_{PQ}D_Q^{-1}B_{QP}B_P^{-1}Pc
   -B_P^{-1}B_{PQ}D_Q^{-1}Qc\bigr]+O(\Gamma^{-2}),\\
x_Q&=(\Gamma D_Q+B_{QQ})^{-1}\bigl(Qc-B_{QP}x_P\bigr)
   =\tfrac1\Gamma D_Q^{-1}\bigl(Qc-B_{QP}B_P^{-1}Pc\bigr)+O(\Gamma^{-2}).
\end{align*}
Contracting with $p^\dagger$ and collecting orders gives
$F_\Gamma=p^\dagger Px_P+p^\dagger Qx_Q
=F_0+\Gamma^{-1}F_1+O(\Gamma^{-2})$ with $F_0$, $F_1$ as in
\eqref{eq:F0}--\eqref{eq:F1}; every remainder constant is uniform on $K$
because all block norms are bounded on $K$ and $\Gamma>\Gamma_*$.
\end{proof}

\begin{remark}[Role of $F_1$]
When $F_0\equiv0$ but the singular structure (H1)--(H4) holds, the class is
decided by the first nonvanishing coefficient among $F_1,F_2,\dots$
(computable by extending the expansion above); if all vanish, Theorem
\ref{thm:I} applied to a realization of $F_\Gamma$ certifies Class I. This is
Step 3 of the decision algorithm.
\end{remark}

\subsection{Insufficiency of endpoint conditions}

\begin{lemma}[$Pc\ne0$ and $p^\dagger P\ne0$ do not suffice]
\label{lem:counterexample}
There exist systems satisfying (H1)--(H4) with $Pc\ne0$ and
$p^\dagger P\ne0$ but $F_0\equiv0$.
\end{lemma}

\begin{proof}[Proof (explicit example)]
On $P\mathcal V\cong\mathbb C^2$ take
\[
B_P=\begin{pmatrix}0&1\\ 1&0\end{pmatrix}
\quad(\text{so } B_P^{-1}=B_P),\qquad
Pc=\begin{pmatrix}1\\0\end{pmatrix},\qquad
P^\dagger p=\begin{pmatrix}1\\0\end{pmatrix}.
\]
Then $F_0=p^\dagger PB_P^{-1}Pc=(1\;0)\,B_P\,(1\;0)^T=0$ although both
endpoint projections are nonzero. Embedding this block into any system with
invertible $D_Q$ (e.g.\ $D=0\oplus I$) satisfies (H1)--(H4). If $B_P$ carries
a $z$-dependent detuning on its diagonal the cancellation can even hold
identically in $z$ whenever the diagonal matrix elements
$(B_P)_{11}$ vanish identically, e.g.\ by a selection rule.
\end{proof}

\begin{remark}
Likewise, leakage $PBQ\ne0$ from the protected subspace into the damped sector
is \emph{permitted}: it only renormalizes $F_1$ and higher orders
(cf.\ \eqref{eq:F1}), not the leading term $F_0$. Protection is a statement
about the leading order, not about perfect isolation.
\end{remark}

\subsection{EIT sector-cut corollary}

\begin{corollary}[Protected EIT channel]
\label{cor:eit}
Let the full and sector-cut generators share the singular structure
(H1)--(H2) with common $P$, and let both protected blocks
$B_{P,\mathrm{full}}(z)$ and $B_{P,\mathrm{cut}}(z)$ satisfy (H3) on $K$.
Then
\[
\delta\chi_S(z)\;=\;\delta\chi_{S,0}(z)+O(\Gamma^{-1}),
\qquad
\delta\chi_{S,0}(z)
=p^\dagger P\bigl[B_{P,\mathrm{full}}(z)^{-1}
 -B_{P,\mathrm{cut}}(z)^{-1}\bigr]Pc ,
\]
uniformly on $K$. If $\delta\chi_{S,0}\not\equiv0$, the sector-mediated
response survives at $O(1)$ under arbitrarily strong dissipation
(Class III for the sector channel); if $\delta\chi_{S,0}\equiv0$, the sector
channel is decided at subleading order as in the remark above.
\end{corollary}

\begin{proof}
Apply Theorem \ref{thm:III} separately to
$\chi_{\mathrm{full}}=p^\dagger[\Gamma D+B_{\mathrm{full}}]^{-1}c$ and
$\chi_{\mathrm{cut}}=p^\dagger[\Gamma D+B_{\mathrm{cut}}]^{-1}c$ (the sector
cut modifies couplings inside $B$, not the dissipation scale structure, by the
construction of companion Theorem 2B) and subtract the two uniform expansions.
\end{proof}

% ==================================================================
\section{Unification: the index $\nu$ and the exclusive trichotomy}
% ==================================================================

\begin{definition}[Response class index]
\label{def:nu}
For a model on the window $K$ define
\[
\nu:=\sup\Bigl\{s\ge0:\ \sup_{z\in K}\Gamma^{\,s}\,|F_\Gamma(z)|
\xrightarrow[\Gamma\to\infty]{}0\Bigr\}
\in[0,\infty],
\]
with the convention $\nu=\infty$ when the supremand vanishes identically for
large $\Gamma$.
\end{definition}

\begin{theorem}[Unified trichotomy]
\label{thm:trichotomy}
Assume the hypotheses of exactly one of the following regimes hold on $K$:
(a) Theorem \ref{thm:I} conditions for a realization of $F_\Gamma$
(all moments zero); (b) Theorem \ref{thm:II} with first nonzero moment
$\mu_m\not\equiv0$; (c) Theorem \ref{thm:III} with $F_0\not\equiv0$.
Then respectively
\[
\text{(a) } \nu=\infty\ \ (\text{Class I}),\qquad
\text{(b) } \nu=\nudiss=m+1\in(0,\infty)\ \ (\text{Class II}),\qquad
\text{(c) } \nu=0\ \ (\text{Class III}),
\]
and the three classes are mutually exclusive and exhaust all models for which
the decision algorithm of Section \ref{sec:algorithm} terminates with a
verdict.
\end{theorem}

\begin{proof}
\emph{(a).} $F_\Gamma\equiv0$ for all $\Gamma$ off a finite set (Theorem
\ref{thm:I}(iv) with resolvent variable $\Gamma$ or $z$ as appropriate), so
every $s$ qualifies and $\nu=\infty$.

\emph{(b).} By Theorem \ref{thm:II}(ii),
$\sup_K\Gamma^s|F_\Gamma|\le C\,\Gamma^{s-(m+1)}\to0$ for $s<m+1$, so
$\nu\ge m+1$. Conversely pick $z_0$ with $\mu_m(z_0)\ne0$: then
$\Gamma^{m+1}|F_\Gamma(z_0)|\to|\mu_m(z_0)|>0$, so no $s>m+1$ (nor $s=m+1$)
qualifies, giving $\nu=m+1$ exactly.

\emph{(c).} Pick $z_0$ with $F_0(z_0)\ne0$: $|F_\Gamma(z_0)|\to|F_0(z_0)|>0$,
so already $s=0^+$ fails and $\nu=0$.

\emph{Exclusivity.} The three values $\nu=\infty$, $\nu\in(0,\infty)$,
$\nu=0$ are disjoint, and $\nu$ is a single well-defined number for each
model (Definition \ref{def:nu} involves only the invariant function
$F_\Gamma$; basis independence is Lemma \ref{lem:basis}). Hence no model can
carry two class labels. \emph{Exhaustiveness relative to the algorithm}: the
algorithm below returns a verdict only by verifying (a), (b), or (c) (or (b)/(a)
at subleading order in the singular case), so every decided model receives
exactly one class.
\end{proof}

\begin{remark}[Uniqueness of asymptotic coefficients]
The proof uses only that limits of $\Gamma^{s}F_\Gamma$ are unique --- the
Laurent/asymptotic coefficients of $F_\Gamma$ in $\Gamma^{-1}$ are unique
whenever they exist. Two different leading orders for the same model are
therefore impossible, which is the precise content of ``exclusive
classification''.
\end{remark}

% ==================================================================
\section{The finite decision algorithm}
\label{sec:algorithm}
% ==================================================================

\begin{theorem}[Termination and correctness]
\label{thm:algorithm}
Given $(D,B(z),c,p)$ on $\Vresp$ (dimension $N$) and a realization
$(M,r,\ell)$ of the sector-cut response (dimension $n$), the following
procedure terminates after finitely many exact linear-algebra operations and
outputs the correct class of Theorem \ref{thm:trichotomy}:
\begin{enumerate}[label=\textbf{Step \arabic*.},leftmargin=5em,nosep]
\item Compute $\mu_k=\ell^\dagger M^kr$, $k=0,\dots,n-1$. If all vanish
      (identically in $z$ where applicable): \textbf{Class I}, stop.
\item If $D$ is invertible on $\Vresp$: compute
      $\mu_k(z)=p^\dagger(D^{-1}B)^kD^{-1}c$ for $k=0,1,\dots$ until the
      first $m\le N-1$ with $\mu_m\not\equiv0$ (guaranteed to exist by
      Theorem \ref{thm:II}(iii), else Step 1 would have fired):
      \textbf{Class II} with $\nudiss=m+1$, stop.
\item If $D$ is singular: verify semisimplicity of the zero eigenvalue,
      construct the Riesz projection $P$, and evaluate $F_0$
      (or $\delta\chi_{S,0}$ of Corollary \ref{cor:eit}). If
      $F_0\not\equiv0$: \textbf{Class III}, stop. If $F_0\equiv0$: evaluate
      $F_1,F_2,\dots$ from the block expansion of Theorem \ref{thm:III};
      the first nonzero coefficient yields \textbf{Class II} with the
      corresponding index; if the expansion terminates with all coefficients
      zero (certified via Step 1 applied to a realization of $F_\Gamma$):
      \textbf{Class I}.
\end{enumerate}
\end{theorem}

\begin{proof}
\emph{Termination.} Step 1 is $n$ matrix--vector products. Step 2 needs at
most $N$ moments by Theorem \ref{thm:II}(iii). Step 3: semisimplicity and $P$
are a single eigendecomposition; each coefficient $F_j$ is a finite block
expression, and the loop over $j$ is capped: $F_\Gamma$ restricted to the
large-$\Gamma$ regime is a rational function of $\Gamma$ of degree bounded by
$N$ (Cramer's rule on $\Gamma D+B$), so its asymptotic expansion is
determined by its first $2N+1$ coefficients; if $F_0=\dots=F_{2N}\equiv0$ the
rational function is identically zero, and the Krylov certificate of Step 1
(applied to the $\Gamma$-realization with $M'=-D^{-1}B$ on the appropriate
subspace, or directly to the rational function) confirms Class I.

\emph{Correctness.} Each verdict is precisely the verified hypothesis of the
corresponding clause of Theorem \ref{thm:trichotomy}: Step 1 verifies (a) via
Theorem \ref{thm:I}(i)$\Rightarrow$(iv); Step 2 verifies (b) via Theorem
\ref{thm:II}; Step 3 verifies (c) via Theorem \ref{thm:III}, or reduces to
(b)/(a) at subleading order as in the Remark after Theorem \ref{thm:III}.
\end{proof}

% ==================================================================
\section{Scope remarks and out-of-scope behaviour}
% ==================================================================

\begin{remark}[Jordan blocks at zero]
If the zero eigenvalue of $D$ is \emph{not} semisimple, Lemma \ref{lem:riesz}
still yields the spectral splitting, but $DP\ne0$ (nilpotent part), the block
$\Gamma DP$ no longer vanishes, and the expansion of Theorem \ref{thm:III}
acquires mixed orders (in general fractional powers of $\Gamma^{-1}$ from the
perturbed nilpotent block). The classification then requires the Newton
polygon of $\det(\Gamma D+B)$; we exclude this case here and flag it as an
assumption (H1), consistent with the physical origin of $D$ as a sum of
Lindblad dissipators, which in the models used in this repository is
diagonalizable on the relevant blocks.
\end{remark}

\begin{remark}[Declared scope]
All results are statements about finite-dimensional, Markovian, weak-probe
linear response, and the classification claim should be quoted as
``material-agnostic structural classification within finite-dimensional
Markovian weak-probe response''. Nonlinear response, non-Markovian memory,
continua, propagation, gain, and time-dependent control require separate
extensions.
\end{remark}

% ==================================================================
\section*{Dependency chart}
% ==================================================================

\begin{center}
\begin{tabular}{lll}
Lemma \ref{lem:basis}, \ref{lem:kalman} & $\longrightarrow$ &
 all theorems (basis/realization independence)\\[2pt]
Companion Thm 2B (sector cut) & $\longrightarrow$ &
 realization \eqref{eq:realization}; Corollary \ref{cor:eit}\\[2pt]
Cayley--Hamilton + adjugate + Laurent & $\longrightarrow$ &
 Theorem \ref{thm:I}\\[2pt]
Neumann series, uniform on compacts & $\longrightarrow$ &
 Theorem \ref{thm:II}\\[2pt]
Riesz projection (Lemma \ref{lem:riesz}) + Schur complement
 & $\longrightarrow$ & Theorem \ref{thm:III}\\[2pt]
Thms \ref{thm:I}+\ref{thm:II}+\ref{thm:III} & $\longrightarrow$ &
 Trichotomy (Thm \ref{thm:trichotomy}) $\longrightarrow$
 Algorithm (Thm \ref{thm:algorithm})
\end{tabular}
\end{center}

\end{document}
